\chapter{Conclusions}
\label{ch:conclusions}

% --------------------------------------------------------------------------
\section{Summary}
\label{sec:summary}
% --------------------------------------------------------------------------

This thesis set out to answer whether a language model, fine-tuned on real BPF
program--verifier-log pairs and trained with reinforcement learning against live
kernel coverage, could drive systematic exploration of the Linux eBPF verifier.
Four contributions were made toward that goal.

\textbf{End-to-end pipeline.}
A complete research platform was designed and implemented: a modified buzzer
instance dumps the kernel FFI boundary as JSONL; a QLoRA fine-tuning stage
trains on the collected dataset; a GRPO RL stage uses KCOV kernel PC coverage
as the reward; and a Go binary running inside a KCOV-instrumented Linux~6.8 VM
validates each generated program and returns its PC trace.
The central engineering challenge --- bridging LLM text output to the raw 8-byte
BPF instruction structs the kernel requires --- is solved by the pure-Python BPF
encoder, which extracts opcode bytes directly from the verifier-log format and
imposes no clang dependency at inference time.
The pipeline is modular: a different research goal (targeting a different kernel
subsystem, for instance) requires only swapping the SFT training data and the
reward function; the scaffolding remains unchanged.

\textbf{SFT result.}
Qwen2.5-Coder-1.5B, fine-tuned with QLoRA on a balanced dataset of 27,514
BPF program--verifier-log pairs, achieves a verifier pass-rate of 60\%,
compared to 1\% for the zero-shot base model --- a 60-fold improvement.
The fine-tuning ran on an RTX~4070~Laptop GPU with 8.59~GB VRAM,
demonstrating that parameter-efficient fine-tuning (0.14\% of model parameters
trained) makes the approach accessible on consumer hardware.
The result confirms RQ1: a 1.5-billion-parameter code model can acquire
sufficient knowledge of BPF verifier syntax from labeled data.

\textbf{Failure-mode diagnosis.}
RL run~1 (\texttt{rl\_grpo\_v2}, 8,370 steps) produced 137 new kernel PCs in the
first $\approx$1,241 batches, then stalled completely.
Post-run analysis established that the reward standard deviation $\sigma_r$
dropped to zero in 99.5\% of batches, zeroing the GRPO gradient and freezing
the model.
The structural cause was twofold.
First, the SFT model --- trained on a \texttt{Status:~VALID} prompt --- learned
to generate almost exclusively valid programs during RL inference (99.7\% accepted
rate), so once the coverage frontier was saturated all four completions per group
received reward~0.1, producing zero variance.
Second, a code-level bug in \texttt{kcov\_validator} discarded KCOV traces for
rejected programs, meaning that the small minority of rejected completions also
contributed nothing to within-group differentiation.
The comparison run with a neutral prompt (63.2\% rejected rate vs.\ 0.3\% in
run~1) provided direct empirical evidence that the prompt was the dominant
attractor.
The KCOV bug was fixed in commit~\texttt{c8a9d02}; because both confounds coexisted
in run~1, their individual contributions cannot be cleanly separated without a
second run using the neutral prompt.
The GRPO gradient-starvation failure mode, once identified, is both diagnosable
from standard training logs and fixable without hyperparameter search.

\textbf{Reward redesign.}
The depth-based verdict-blind reward removes the accept/reject split from the
reward formula entirely.
Programs are scored by the number of kernel PCs reached during verification and
by a bonus for PCs not seen in prior batches, independent of the verifier verdict.
This design structurally guarantees non-zero within-group reward variance: even
when all $G$ completions are rejected, programs that trigger deeper verifier
analysis receive higher depth scores and therefore different rewards.
The redesigned reward function is implemented in \texttt{ml/reward.py} and
validated in the 100-step comparison run; a full training run (run~2) is the
next step.

Together, the four contributions demonstrate that LLM-guided eBPF verifier
fuzzing with live kernel coverage feedback is a technically feasible research
direction.
The SFT stage resolves the structural-validity bottleneck that defeats random
mutation.
The failure-mode analysis shows that GRPO reward design is the critical variable
for sustained exploration: a reward function that guarantees within-group variance
is not an implementation detail but a correctness requirement for GRPO-based
training.

% --------------------------------------------------------------------------
\section{Future Work}
\label{sec:future-work}
% --------------------------------------------------------------------------

\begin{enumerate}

  \item \textbf{RL run~2.}
        Execute with the depth-based verdict-blind reward
        (\texttt{ml/reward.py}) and a neutral prompt; this isolates the prompt
        bias from the KCOV bug.
        The Colab Pro remote-reward pipeline (\texttt{ml/train\_grpo\_colab.ipynb},
        \texttt{tools/reward\_server.py}) is implemented and ready.

  \item \textbf{SFT toward verifier stress.}
        Retrain with prompts targeting rejected programs or programs that reach
        deep verifier paths rather than valid programs; evaluate whether this
        shifts RL exploration away from the valid-program attractor established
        by the current SFT objective.

  \item \textbf{Curriculum learning.}
        Staged fine-tuning: a syntax acquisition phase (optimise for valid
        programs) followed by a stress phase (optimise for rejected or deep-path
        programs).  This mirrors curriculum learning techniques used in
        mathematical reasoning models.

  \item \textbf{GRPO reward shaping.}
        Apply fixed-reference advantage $A = 2r - 1$ \cite{grpo-starvation}
        to handle binary-reward collapse without relying on within-group
        reward variance.

  \item \textbf{Scale.}
        Apply the same pipeline to a 7B-parameter model on cloud GPU; the
        Colab Pro remote-reward infrastructure is already built for this.

  \item \textbf{Systematic comparison.}
        Buzzer \texttt{coverage\_based} strategy vs.\ LLM-guided generation at
        equal wall-clock time, measured in unique verifier PCs explored.
        This would quantify whether LLM-guided generation provides a coverage
        advantage over the existing state of the art.

  \item \textbf{Constrained decoding.}
        Enforce eBPF instruction semantics at generation time (e.g., via a
        finite-state-machine constraint on the token sequence) to reduce the
        encode-fail rate and increase the density of reward-eligible programs
        per training batch.

\end{enumerate}
